% Part: first-order-logic
% Chapter: tableaux
% Section: quantifier-rules

\documentclass[../../../include/open-logic-section]{subfiles}

\begin{document}

\olfileid{fol}{tab}{qrl}

\olsection{Regels voor kwantoren}

\subsection{Regels voor $\lforall$ (voor alle)}

\begin{defish}
\AxiomC{\sFmla{\True}{\lforall[x][!A(x)]}}
\RightLabel{\TRule{\True}{\forall}}
\UnaryInfC{\sFmla{\True}{!A(t)}}
\DisplayProof
\hfill
\AxiomC{\sFmla{\False}{\lforall[x][!A(x)]}}
\RightLabel{\TRule{\False}{\lforall}}
\UnaryInfC{\sFmla{\False}{!A(a)}}
\DisplayProof
\end{defish}

Bij \TRule{\True}{\lforall} is $t$ een gesloten term, dus een term
zonder variabelen. Bij \TRule{\False}{\lforall} is $a$~!!a{constant}.
Die constante mag nergens voorkomen in de tak boven de regel
\TRule{\False}{\lforall}. We noemen $a$ de \emph{eigenvariabele} van
de inferentie \TRule{\False}{\forall}.\footnote{We noemen $a$ hier om
historische redenen een ``eigenvariabele'', hoewel hij in de
bovenstaande regel !!a{constant} is.}

\subsection{Regels voor $\lexists$ (er bestaat)}

\begin{defish}
\AxiomC{\sFmla{\True}{\lexists[x][!A(x)]}}
\RightLabel{\TRule{\True}{\lexists}}
\UnaryInfC{\sFmla{\True}{!A(a)}}
\DisplayProof
\hfill
\AxiomC{\sFmla{\False}{\lexists[x][!A(x)]}}
\RightLabel{\TRule{\False}{\lexists}}
\UnaryInfC{\sFmla{\False}{!A(t)}}
\DisplayProof
\end{defish}

Ook hier is $t$~een gesloten term. Verder is $a$~!!a{constant}. Deze
constante mag niet voorkomen in de tak boven de
regel~\TRule{\True}{\lexists}. We noemen $a$ de
\emph{eigenvariabele} van de inferentie
\TRule{\True}{\lexists}.

Bij deze twee regels geldt dus dezelfde eis: de eigenvariabele mag niet
voorkomen in de tak boven de inferentie \TRule{\False}{\lforall} of
\TRule{\True}{\lexists}. Die eis heet de
\emph{eigenvariabelevoorwaarde}.

\begin{explain}
We gebruiken de volgende afspraak. Als $!A$ !!a{formula} is waarin de
!!{variable}~$x$ vrij voorkomt, schrijven we dat als~$!A(x)$. In
dezelfde situatie is $!A(t)$ een korte schrijfwijze voor
$\Subst{!A}{t}{x}$. Daarom kunnen we de regel
\TRule{\False}{\lexists} ook als volgt schrijven:
\begin{prooftree}
  \AxiomC{\sFmla{\False}{\lexists[x][!A]}}
  \RightLabel{\TRule{\False}{\lexists}}
  \UnaryInfC{\sFmla{\False}{\Subst{!A}{t}{x}}}
\end{prooftree}
De term $t$ kan al in~$!A$ staan. Zo kan $!A$ bijvoorbeeld
$\Atom{\Obj P}{t,x}$ zijn. Daarom mag je
$\sFmla{\False}{\Atom{\Obj P}{t,t}}$ afleiden uit
$ \sFmla{\False}{\lexists[x][\Atom{\Obj P}{t,x}]}$ met
\TRule{\False}{\lexists}. Bij \TRule{\False}{\lforall}
en~\TRule{\True}{\lexists} geldt een extra voorwaarde: de
!!{constant}~$a$ mag niet in~$!A$ voorkomen. Daarom mag je
$\sFmla{\False}{\Atom{\Obj P}{a,a}}$ niet afleiden uit
$\sFmla{\False}{\lforall[x][\Atom{\Obj P}{a,x}]}$ met
$\TRule{\False}{\lforall}$.
\end{explain}

\begin{explain}
Bij \TRule{\True}{\lforall} en \TRule{\False}{\lexists} mag $t$ elke
term zijn. Bij \TRule{\True}{\lexists} en
\TRule{\False}{\lforall} geldt wel de eigenvariabelevoorwaarde. De
!!{constant}~$a$ mag dan nergens voorkomen in de takken boven de
betreffende inferentie. Deze voorwaarde zorgt ervoor dat het systeem
alleen geldige conclusies oplevert. Zonder de voorwaarde zou het
volgende een gesloten !!{tableau} zijn voor
$\lexists[x][\formula{A}(x)] \lif \lforall[x][\formula{A}(x)]$:
\begin{center}
\begin{tableau}{}
  [\sFmla{\False}{\lexists[x][\formula{A}(x)] \lif \lforall[x][\formula{A}(x)]}, just=\TAss
    [\sFmla{\True}{\lexists[x][\formula{A}(x)]},
      just={\TRule{\False}{\lif}[1]}
      [\sFmla{\False}{\lforall[x][\formula{A}(x)]},
        just={\TRule{\False}{\lif}[1]}
        [\sFmla{\True}{\formula{A}(a)},
          just={\TRule{\True}{\lexists}[2]}
          [\sFmla{\False}{\formula{A}(a)},
            just={\TRule{\False}{\lforall}[3]}, close]
        ]
      ]
    ]
  ]
\end{tableau}
\end{center}
Maar $\lexists[x][\formula{A}(x)] \lif
\lforall[x][\formula{A}(x)]$ is niet geldig.
\end{explain}

\end{document}
